\section{Method}
\label{sec:method}

\subsection{Problem Definition}

Each sample is anchored at a label time $\tau$ and is assigned one of eleven machinery operation classes. The model receives a synchronized multimodal context
\begin{equation}
    \vx = \{\vx^{t}, \vx^{v}, \vx^{a}\},
\end{equation}
where $\vx^{t} \in \R^{T \times F_t}$ is a causal trajectory feature sequence, $\vx^{v}$ is a set of ViT visual frames selected from the same causal window, and $\vx^{a}$ is the audio segment aligned to the anchor time. The task is to learn
\begin{equation}
    \hat{y} = \argmax_{c \in \{0,\ldots,10\}} p(y=c \mid \vx).
\end{equation}
Throughout this paper, the task is referred to as \task{}.

\subsection{Temporal Alignment}

All modalities are aligned around the same anchor time. The trajectory branch uses a causal window of $T=128$ seconds. This window contains only observations at or before $\tau$, preventing future information leakage. The 128-second length is selected from the completed seq\_len ablation in \Cref{sec:seq_ablation}.

The ViT branch follows the same 128-second causal window. Rather than uniformly sampling frames across the whole window, we select the nine frames closest to $\tau$ among frames whose timestamps are not later than $\tau$. This preserves visual evidence near the label while respecting causal evaluation. If fewer than nine frames are available, the nearest valid frames are reused by the data loader.

The audio branch uses the synchronized audio clip associated with the anchor time. In the current implementation, audio is transformed into a log-mel spectrogram and encoded by an AST backbone.

\begin{figure}[t]
\centering
\includegraphics[width=0.98\linewidth]{fig_method_data_processing.png}
\caption{Data processing and temporal alignment workflow. Each label is associated with one anchor time $\tau$. The trajectory branch receives a 128-s causal sequence, the ViT branch receives the nine nearest causal frames from the same window, and the audio branch receives the one-second segment aligned to the anchor. The three processed streams form one matched multimodal sample for classification.}
\label{fig:method_data_processing}
\end{figure}
\FloatBarrier

\subsection{Trajectory Encoder}
\label{sec:trajectory_encoder}

The trajectory branch describes the continuous motion state of the machine before the label time. For each anchor time $\tau$, the input is a causal sequence of $T=128$ one-second samples. A raw trajectory point is denoted as
\begin{equation}
    X_i = (\lambda_i, \phi_i, v_i, d_i, \theta_i, t_i),
\end{equation}
where $\lambda_i$ and $\phi_i$ are longitude and latitude, $v_i$ is speed, $d_i$ is the recorded cutting-depth-related variable, $\theta_i$ is the heading angle, and $t_i$ is the timestamp. These original attributes provide the instantaneous measurement at one point, but they do not fully describe how the machine is moving over a short period. This limitation is important for \task{}, because several operation classes differ mainly in motion evolution. For example, straight harvesting and transfer can share similar speeds, while turning empty harvesting and turning transfer require the model to compare heading, curvature, speed, and working-context changes over time.

We therefore use a multi-view trajectory feature extraction procedure before the neural encoder. It has two causal stages: instantaneous physical feature extraction (IPE) and multi-window statistical extraction (MWSE). IPE converts the original GNSS attributes into interpretable kinematic variables. To avoid artificial discontinuities around $0^\circ/360^\circ$, heading differences are computed with circular wrap-around:
\begin{equation}
    \Delta^{c}\theta_i = ((\theta_i - \theta_{i-1} + 180) \bmod 360) - 180 .
\end{equation}
Let $\Delta t_i=t_i-t_{i-1}$. Since the valid trajectory stream is sampled at 1 Hz inside each contiguous segment, $\Delta t_i=1$ s in normal samples; the formulas below are written in rate form for clarity. The core IPE variables are
\begin{align}
    a_i &= \frac{v_i-v_{i-1}}{\Delta t_i}, \\
    s_i &= v_i-v_{i-1}, \\
    \omega_i &= \frac{\Delta^{c}\theta_i}{\Delta t_i}, \\
    \alpha_i &= \frac{\omega_i-\omega_{i-1}}{\Delta t_i}, \\
    g_i &= \omega_i-\omega_{i-1},
\end{align}
where $a_i$ is acceleration, $s_i$ is speed difference, $\omega_i$ is angular speed, $\alpha_i$ is angular acceleration, and $g_i$ is angular-speed difference. In the implementation, these quantities are computed separately within each contiguous video segment; the first sample of a segment uses zero first-order difference by prepending itself.

Additional IPE descriptors are included to represent spatial displacement, work-state variation, and turning intensity:
\begin{align}
    \Delta \lambda_i &= \lambda_i-\lambda_{i-1}, &
    \Delta \phi_i &= \phi_i-\phi_{i-1}, &
    \Delta d_i &= d_i-d_{i-1}, \\
    j_i &= a_i-a_{i-1}, &
    \kappa_i &= \frac{\omega_i}{|v_i| + 10^{-3}}, &
    e_i &= v_i^2 + |a_i|.
\end{align}
Here $j_i$ is jerk, $\kappa_i$ is a speed-normalized curvature proxy, and $e_i$ is a motion-energy descriptor. The absolute heading change $|\Delta^{c}\theta_i|$ is used as a turning indicator. These instantaneous variables make the trajectory representation sensitive to acceleration, heading transition, local curvature, depth variation, and stop-and-go behavior rather than relying only on raw coordinates and speed.

MWSE then summarizes the recent trajectory context using causal rolling windows. For a scalar IPE series $q_i$, the causal window of length $w$ is
\begin{equation}
    \mathcal{H}_w(i)=\{q_k \mid \max(1,i-w+1) \leq k \leq i\}.
\end{equation}
For $q_i \in \{v_i,a_i,\omega_i,\alpha_i,\Delta^{c}\theta_i\}$ and $w \in \{5,20\}$ seconds, we compute the rolling median and standard deviation:
\begin{equation}
    \operatorname{MWSE}_{5,20}(q_i)=
    [\operatorname{med}(\mathcal{H}_5(i)), \operatorname{med}(\mathcal{H}_{20}(i)),
    \operatorname{std}(\mathcal{H}_5(i)), \operatorname{std}(\mathcal{H}_{20}(i))].
\end{equation}
The 5-s window captures short maneuver changes, while the 20-s window captures more stable operation patterns. We also compute a 3-s local speed mean $\bar{v}^{(3)}_i$ and a 3-s local heading standard deviation $\sigma^{(3)}_{\theta,i}$ from the current and two previous samples. All IPE and MWSE features are causal and therefore use only observations at or before the current time.

The final trajectory feature vector has 36 dimensions:
\begin{equation}
\begin{split}
    \vx_i^t = [&\lambda_i, \phi_i, v_i, a_i, \omega_i, \alpha_i, \Delta^{c}\theta_i, \\
    &\operatorname{MWSE}_{5,20}(v_i,a_i,\omega_i,\alpha_i,\Delta^{c}\theta_i), \\
    &\Delta \lambda_i, \Delta \phi_i, j_i, \kappa_i, e_i, |\Delta^{c}\theta_i|,
    \bar{v}^{(3)}_i, \sigma^{(3)}_{\theta,i}, \Delta d_i].
\end{split}
\end{equation}
The feature construction is performed before sample extraction and is restricted to each contiguous video segment. When a causal 128-s sample extends before the segment beginning, the earliest valid sample is repeated; no observations from other segments or from the future are used. The resulting $\vx^{t}\in\R^{128\times36}$ sequence is normalized using the mean and standard deviation estimated from the training split only.

After feature extraction, the sequence is encoded by a BiLSTM-style neural encoder. Each 36-D vector is reshaped into a $6\times6$ feature map:
\begin{equation}
    M_i=\operatorname{reshape}(\vx_i^t,6,6).
\end{equation}
A point-wise CNN with two convolution--batch-normalization--ReLU blocks and max pooling extracts local interactions among the kinematic variables:
\begin{equation}
    \vp_i = \psi_{\mathrm{cnn}}(M_i).
\end{equation}
The projected point representations $\{\vp_i\}_{i=1}^{128}$ are then passed to a two-layer bidirectional LSTM. The forward and backward states inside the causal window are summed at each time step,
\begin{equation}
    \vh_i = \overrightarrow{\vh}_i + \overleftarrow{\vh}_i,
\end{equation}
and an attention pooling layer aggregates the 128 hidden states into a fixed-dimensional trajectory embedding:
\begin{align}
    \beta_i &= \frac{\exp(\mathbf{w}^{\top}\tanh(\mathbf{W}\vh_i))}
    {\sum_{k=1}^{128}\exp(\mathbf{w}^{\top}\tanh(\mathbf{W}\vh_k))}, \\
    \vz^{t} &= \sum_{i=1}^{128}\beta_i \vh_i = f_t(\vx^{t}).
\end{align}
Although the LSTM is bidirectional within the input sequence, the whole sequence is causal with respect to $\tau$, so the trajectory encoder never uses observations after the label time. This design separates interpretable point-level kinematic construction from sequence-level temporal reasoning. It is especially useful for turning-related and transfer-related behaviors, where the decisive evidence is the recent evolution of heading, curvature, speed, and depth-related variables rather than any single GNSS point.

\subsection{ViT Visual Encoder}

The ViT modality captures visual context such as field texture, crop intake, road surface, and the machine's local working environment. Each selected frame is resized to $224 \times 224$ and encoded by a pretrained ViT-style visual backbone~\cite{dosovitskiy2021vit}. The current best ViT configuration uses nine nearest causal frames, temporal GRU pooling, and frame-difference features. For frame embeddings $\{\vh_i^{v}\}_{i=1}^{9}$, the temporal pooling module produces
\begin{equation}
    \vz^{v} = f_v(\vh_1^{v}, \ldots, \vh_9^{v}).
\end{equation}
This design keeps the visual evidence close to the decision time and avoids using visual frames from the future.

\subsection{Audio Encoder}

The audio modality captures engine load, idling, unloading, and mechanical-component engagement. Audio clips are transformed into log-mel spectrograms and encoded using an Audio Spectrogram Transformer initialized from AudioSet-pretrained weights~\cite{gong2021ast}. The audio encoder outputs
\begin{equation}
    \vz^{a} = f_a(\vx^{a}).
\end{equation}
Audio is expected to be especially useful for distinguishing classes with similar or stationary trajectories, such as engine-off waiting, idling waiting, and unloading.

\subsection{Class-Adaptive Trimodal Fusion}
\label{sec:class_gate}

The proposed trimodal model uses class-adaptive logit fusion instead of direct feature concatenation. This design is motivated by the observation that different operation classes are dominated by different evidence sources. Motion-dependent classes require trajectory curvature, heading change, and speed transitions; scene-dependent classes benefit from ViT visual context; process states such as engine-off waiting, idling waiting, and unloading are often better captured by audio.

For modality $m \in \{t,v,a\}$, a modality-specific classifier head maps the embedding $\vz^m$ to eleven-class logits:
\begin{equation}
    \vl^{m} = q_m(\vz^m), \quad \vl^{m} \in \R^{11}.
\end{equation}
A gating network receives the same three modality embeddings and outputs a class-specific softmax weight $g_{c}^{m}$ over modalities:
\begin{equation}
    \sum_m g_{c}^{m}=1, \quad g_{c}^{m} \geq 0 .
\end{equation}
The final logit for class $c$ is
\begin{equation}
    l_c = \sum_m g_{c}^{m} l_c^{m}.
\end{equation}
The predicted distribution is then
\begin{equation}
    p(y \mid \vx) = \softmax(\vl).
\end{equation}

In the implementation, each modality first passes through layer normalization, dropout, and a lightweight classifier head. The gate projects each modality embedding into a shared gate space, concatenates the three gate features, and predicts an $11 \times 3$ matrix of class-modality weights. The final gate layer is initialized to zero, so the model starts from nearly uniform modality weights and learns class-specific preferences during training. This formulation allows the model to preserve audio-dominant evidence for process classes while assigning more weight to BiLSTM or ViT evidence for motion and scene classes.

The trimodal model is initialized from the unimodal checkpoints for the same seed. The encoders are frozen for the first epoch and then jointly trained, which stabilizes the fusion module while preserving the unimodal representations learned in the preceding runs.

\subsection{Concatenation Reference}

For reference, we also keep a direct concatenation baseline:
\begin{equation}
    \vz = [\vz^{t}; \vz^{v}; \vz^{a}],
\end{equation}
followed by layer normalization, dropout, and a multilayer classifier. This baseline establishes whether the three modalities are complementary under a matched temporal protocol, but it uses one shared classifier for all classes and therefore cannot explicitly preserve class-specific modality dominance.

\subsection{Training Objective and Model Selection}

All models are trained with class-weighted cross-entropy to reduce the dominance of frequent operation classes:
\begin{equation}
    \cL = -\frac{1}{N}\sum_{i=1}^{N} w_{y_i}\log p(y_i \mid \vx_i),
\end{equation}
where $w_{y_i}$ is derived from the inverse class frequency with a power coefficient. The main selection metric is validation macro-F1:
\begin{equation}
    \mathrm{MacroF1} = \frac{1}{11}\sum_{c=0}^{10} \mathrm{F1}_{c}.
\end{equation}
Macro-F1 weights every class equally and is therefore more appropriate than accuracy for this imbalanced eleven-class problem. A checkpoint is saved only when validation macro-F1 improves by more than the early-stopping threshold.
